\begin{afterword}
\summaryhead

The essay opens with a passage from Robyn Dawes about Earl Warren's testimony in February 1942. Told that Japanese-Americans had committed no sabotage, Warren replied that this absence was ``the most ominous sign in our whole situation.'' The author analyzes the reply in probability terms. Evidence favors the hypotheses that made it more likely. Whatever the chance that a Fifth Column would commit no sabotage, the chance is higher if there is no Fifth Column, so the absence of sabotage makes a Fifth Column less likely, not more.

The author grants that absence of proof is not proof of absence, but says that in probability theory, if seeing something would raise the probability of a hypothesis, not seeing it must lower it. The effect can be strong or very weak, depending on how reliably the cause would produce the sign. Gaps in the fossil record are weak evidence; a complete lack of positive observations is strong. A model is tested by what it forbids.

\responsehead

The essay takes what its own source calls ``a most grievous chapter in United States history'' and uses it as an exercise in likelihood ratios. It gets the exercise right in outline. It does not notice what the exercise leaves out.

Warren's error, on the essay's analysis, was to read the absence of sabotage in the wrong direction. That is true and it is not what was wrong with internment. Warren himself, in his memoir as quoted on Wikipedia, regretted the removal ``because it was not in keeping with our American concept of freedom and the rights of citizens.'' The essay's frame has no place for that. By the essay's own arithmetic, a few acts of sabotage would have raised the probability of H1. The essay never says that sabotage by a few would still not have justified imprisoning a population, and a reader who learns only its lesson is not told. The essay treats the case purely as arithmetic.

The case's facts are also wrong as given. Warren was Attorney General, not governor, in February 1942. The post passes on its source's error.

The title attacks a slogan it never quotes or attributes. In Sagan's use, ``absence of evidence is not evidence of absence'' is a warning against treating what has not been disproved as proved. The essay agrees with that in so many words (``Absence of proof is not proof of absence'') and then declares the slogan wrong by reading ``evidence'' in its own technical sense. What remains, once the essay qualifies itself, is that absence may be ``strong evidence of absence or very weak evidence of absence.'' That qualification is the useful rule, and it appears only in the second-to-last paragraph. In ``When Science Can't Help'' (May 2008) the author put the qualification first: absence is evidence ``only to the degree that we could reasonably expect the evidence to appear.'' In this post, readers see the unqualified title first.

In short: a slogan ``refuted'' by changing what its words mean, illustrated with a wartime injustice that the essay treats as a mistake in arithmetic, taken from a source whose facts it did not check.
\end{afterword}
